\section{Ce qui peut mal tourner}\label{sec:risques}
% Chapitre 7 (tache 10, passage chapitre 07). Figures: fig_shiller_div, fig_einbrueche,
% fig_rueckspiel (colonnes ans_tenu/ans_echec/ans_non_juge ajoutees par kapitel_7),
% fig_mc_faecher (colonnes seuil_mc/capital_requis ajoutees par kapitel_7), fig_mc_erfolg,
% fig_inflation_2022. Tableau tab:projection-contre-simulation ecrit ici. Aucun chiffre
% tape: macros de data/kennzahlen.tex (kapitel_7 et _macros_redaction_kapitel_7 dans
% scripts/rechnung_retraite.py). Le seuil de revenu, la longueur des blocs et le nombre
% de trajectoires sont lus dans scripts/montecarlo.py et scripts/histoire.py.
% fig_puffer (reserve de 0 a 3 ans) appartient au chapitre 8: ce chapitre ne cite que
% \ErfolgsquoteBasis et \ErfolgsquoteMitPuffer, issus du MEME calcul.

Ce chapitre r\'epond \`a la question~: que devient ton plan si les dividendes baissent, si les
march\'es ne livrent pas leur moyenne dans l'ordre esp\'er\'e ou si les prix montent plus vite que
pr\'evu~? Le chapitre~\ref{sec:portefeuille} s'est conclu sur un ETF maison qui diversifie l'essentiel
du risque propre \`a chaque titre, mais pas le risque commun \`a tout le march\'e. C'est ce risque
commun que ce chapitre mesure, en rempla\c{c}ant le rendement constant des chapitres~\ref{sec:capital}
et~\ref{sec:accumulation} par l'histoire r\'eelle. La section~\ref{sec:risques-histoire} suit le
dividende r\'eel du march\'e am\'ericain depuis \ShillerDebut{} et le compare \`a l'hypoth\`ese de
croissance du plan, et la section~\ref{sec:risques-baisses} date ses baisses pendant les crises. La
section~\ref{sec:risques-modeles} d\'ecrit le mod\`ele simplifi\'e qu'emploient les sections
suivantes, qu'il ne faut pas confondre avec la projection des chapitres pr\'ec\'edents. Le rejeu de
l'histoire ann\'ee par ann\'ee (section~\ref{sec:risques-rejeu}), le tirage Monte Carlo
(section~\ref{sec:risques-montecarlo}) et le choc de la premi\`ere ann\'ee de retraite
(section~\ref{sec:risques-sequence}) mesurent ensuite la tenue du revenu, puis la
section~\ref{sec:risques-reserve} chiffre l'effet d'une r\'eserve de liquidit\'es. La
section~\ref{sec:risques-inflation} traite enfin de l'inflation de 2022.

\subsection{Le dividende r\'eel du march\'e am\'ericain depuis \ShillerDebut}\label{sec:risques-histoire}

La seule s\'erie publique de dividendes qui couvre plus d'un si\`ecle est celle de l'indice
\IndiceActionsUs{}, tenue par Robert Shiller. La figure~\ref{fig:dividende-reel-indice} en donne le
dividende annuel, corrig\'e de l'inflation.

\annahme{Donn\'ees mensuelles de Robert Shiller\QH{shiller_url}~: cours, dividende et indice des prix
am\'ericains de l'\IndiceActionsUs{}, de \ShillerDebut{} \`a \ShillerFin. Le dividende annuel est la
moyenne des taux annualis\'es publi\'es chaque mois, ramen\'ee aux prix de \ShillerFin{} avec l'indice
des prix am\'ericain. \textbf{Limite d\'eclar\'ee~: march\'e am\'ericain seul.} Aucune s\'erie aussi
longue n'existe dans ce document pour l'Allemagne, la France ou ton ETF maison, et tout ce chapitre en
h\'erite.}

\linienfigur{fig_shiller_div}[xlabel={Ann\'ee},
  ylabel={Dividende r\'eel (dollars de \ShillerFin{} par unit\'e d'indice)},
  xticklabel style={/pgf/number format/set thousands separator={}}, ymin=0,
  legend pos=north west, legend cell align=left, legend style={fill=white}]%
  {jahr}{div_real}%
  {{Dividende annuel de l'\IndiceActionsUs}}%
  {Dividende annuel r\'eel de l'indice \IndiceActionsUs{} de \ShillerDebut{} \`a \ShillerFin, en
  dollars de \ShillerFin{} par unit\'e d'indice (donn\'ees de Robert Shiller, march\'e am\'ericain
  seul).\label{fig:dividende-reel-indice}}

Ce graphique montre que le dividende r\'eel de l'indice a \'et\'e multipli\'e par \ShillerDivMultiple{}
entre \ShillerDebut{} et \ShillerFin, soit une croissance de \ShillerDivCroissance~\% par an au-del\`a
de l'inflation. Cette progression n'a rien de r\'egulier~: de longs paliers et plusieurs chutes
s'intercalent entre les phases de hausse, et la section~\ref{sec:risques-baisses} les date. Sur la
derni\`ere ann\'ee de la s\'erie, \ShillerFin, le dividende r\'eel a encore progress\'e de
\ShillerDivCroissanceDerniere~\% malgr\'e la pouss\'ee d'inflation am\'ericaine.

Cette croissance historique est \`a comparer \`a celle que suppose le plan. Le mod\`ele des
chapitres~\ref{sec:capital} \`a~\ref{sec:portefeuille} fait cro\^itre le dividende de l'ETF maison de
\CroissanceModeleMaison~\% par an en termes r\'eels, ce qui garde son rendement total au niveau du
march\'e, \RenditeReal~\%. Or, sur la fen\^etre \FenetreMarcheDebut--\FenetreMarcheFin{} qui fixe ce
rendement, le dividende de l'indice n'a cr\^u que de \ShillerDivCroissanceFenetre~\% par an. Le reste du
rendement est venu d'une hausse des cours plus rapide que celle des dividendes~: le rendement du
dividende de l'indice est pass\'e de \ShillerRendementDebutFenetre~\% en \FenetreMarcheDebut{} \`a
\ShillerRendementFinFenetre~\% en \FenetreMarcheFin. \`A valorisation constante, le rendement total
d'une action vaut \`a peu pr\`es son rendement du dividende plus la croissance de ce dividende~; avec
\RenditeDivMaison~\% de rendement et la croissance historique de \ShillerDivCroissanceFenetre~\%, on
reste loin des \RenditeReal~\% du plan. \bk{Le plan suppose donc que ton dividende croisse
\CroissanceModeleRapport{} fois plus vite que celui du march\'e am\'ericain depuis
\FenetreMarcheDebut.} C'est une hypoth\`ese favorable~: si ton dividende ne suit que le rythme pass\'e
de l'indice, l'objectif s'\'eloigne, et les \^ages du chapitre~\ref{sec:accumulation} reculent.

\subsection{Les baisses pendant les crises}\label{sec:risques-baisses}

Un revenu de dividendes ne s'effondre pas en un jour comme le cours d'une action lors d'un krach, mais
il peut baisser fortement et longtemps. La figure~\ref{fig:baisses-dividende-crises} recense les
baisses du dividende r\'eel de l'indice.

\annahme{Une baisse est compt\'ee d\`es que le dividende r\'eel annuel recule d'au moins
\SeuilBaisseHistorique~\% sous son dernier sommet~; l'\'episode va de l'ann\'ee du sommet \`a celle du
creux. La baisse est r\'eelle~: elle peut venir d'une coupe du dividende en dollars comme d'une
inflation plus rapide que lui.}

\balkenfigur{fig_einbrueche}[ylabel={Baisse du dividende r\'eel depuis le sommet (\%)},
  ymax=0, enlarge y limits={lower, value=0.2}, bar width=18pt,
  nodes near coords, every node near coord/.append style={font=\scriptsize,
    /pgf/number format/fixed, /pgf/number format/fixed zerofill, /pgf/number format/precision=1}]%
  {episode}{rueckgang_prozent}%
  {Baisses du dividende r\'eel de l'indice \IndiceActionsUs{} d'au moins \SeuilBaisseHistorique~\%
  sous le sommet pr\'ec\'edent, de \ShillerDebut{} \`a \ShillerFin, par \'episode (ann\'ee du sommet et
  ann\'ee du creux).\label{fig:baisses-dividende-crises}}

Ce graphique montre \NombreBaisses{} \'episodes de baisse en un si\`ecle et demi. La pire,
\GroessterEinbruch~\%, court de \GroessterEinbruchDebut{} \`a \GroessterEinbruchJahr, autour de la
Premi\`ere Guerre mondiale~; la suivante, \BaisseLongue~\%, de \BaisseLongueDebut{} \`a
\BaisseLongueFin, couvre la Grande D\'epression et la Seconde Guerre mondiale. La crise financi\`ere a
co\^ut\'e \EinbruchRecent~\% de \EinbruchRecentDebut{} \`a \EinbruchRecentFin. La d\'ecennie
d'inflation, de \EinbruchInflationDebut{} \`a \EinbruchInflationFin, a co\^ut\'e \EinbruchInflation~\%,
surtout parce que les prix montaient plus vite que les dividendes.

La profondeur ne dit pas tout~; la dur\'ee compte autant pour qui vit de ce revenu. Apr\`es le sommet
de \GroessterEinbruchDebut, le dividende r\'eel n'a retrouv\'e son niveau qu'en
\GroessterEinbruchRetour, \GroessterEinbruchSousSommetAns~ans plus tard. Apr\`es celui de
\BaisseLongueDebut, il a fallu \BaisseLongueAns~ans, jusqu'en \BaisseLongueRetour. La baisse de
\EinbruchRecentDebut{} a \'et\'e rattrap\'ee en \EinbruchRecentRetour. Une partie de la presse
sp\'ecialis\'ee pr\'esente les entreprises qui versent des dividendes comme plus stables en p\'eriode
de crise\QP{dividende_statt_rente}~; l'histoire de l'indice montre que leur dividende, lui, n'est pas
\`a l'abri. \bk{Ces baisses sont celles d'un large indice~: un ETF maison de \NombreTitres{} titres peut
faire pire,} comme l'a montr\'e la suppression du dividende de \BeispielBaisse{} au
chapitre~\ref{sec:portefeuille}.

\subsection{Deux mod\`eles \`a ne pas confondre}\label{sec:risques-modeles}

Les sections suivantes emploient un autre mod\`ele que la projection des
chapitres~\ref{sec:capital} et~\ref{sec:accumulation}. Ce mod\`ele est volontairement simple~: il
rejoue l'histoire du march\'e am\'ericain avec ton plan d'\'epargne, mais sans imp\^ot ni assurance
maladie, avec un rendement du dividende fixe et une r\`egle de r\'eussite qui lui est propre. Le
tableau~\ref{tab:projection-contre-simulation} met les deux mod\`eles c\^ote \`a c\^ote.

\begin{table}[!tb]
  \centering\small
  \caption{Ce qui distingue la projection des chapitres~\ref{sec:capital} et~\ref{sec:accumulation}
  du rejeu historique et du Monte Carlo de ce chapitre (euros de 2026).}\label{tab:projection-contre-simulation}
  \begin{tabularx}{\linewidth}{>{\bfseries\raggedright\arraybackslash}p{0.2\linewidth}
    *{2}{>{\raggedright\arraybackslash}X}}
    \toprule
    \rowcolor[gray]{0.9}
    \textbf{Param\`etre} & \textbf{Projection (chapitres~\ref{sec:capital} et~\ref{sec:accumulation})}
      & \textbf{Rejeu et Monte Carlo (ce chapitre)} \\
    \midrule
    Rendement du march\'e & Constant, \RenditeReal~\% r\'eel par an (\FenetreMarcheDebut--\FenetreMarcheFin)
      & Rendements r\'eels annuels de l'\IndiceActionsUs, \ShillerDebut--\ShillerFin{} (moyenne compos\'ee
      \McRenditeHistorique~\%)~: dans l'ordre historique pour le rejeu, par blocs de \McBloc~ans tir\'es
      au hasard pour le Monte Carlo \\
    Imp\^ot et assurance maladie & R\`egles de 2026 (chapitre~\ref{sec:fiscalite})
      & Aucun~: l'objectif est compar\'e au dividende brut \\
    Rendement du dividende & \RenditeDivMaison~\% pour l'ETF maison et \RenditeDivEtf~\% pour l'ETF
      monde, selon la r\'epartition \ReferenzStrategie
      & \RenditeDivMaison~\% sur tout le capital, fixe \\
    \'Epargne & \ReferenzSparquote~\% du salaire net, capital de d\'epart de \KapitalDepartMitte~€
      & M\^eme plan d'\'epargne, sans capital de d\'epart, arr\^et\'e d\`es l'objectif \\
    Objectif atteint & Revenu net d'au moins \ZielMonat~€ par mois
      & Capital fois \RenditeDivMaison~\% d'au moins \ZielJahr~€, soit un capital de \McKapitalSeuil~€ \\
    Revenu \`a la retraite & Non simul\'e~: le capital du chapitre~\ref{sec:capital} suppose un dividende
      qui garde sa valeur r\'eelle
      & Part de l'objectif exact, puis suit chaque ann\'ee la croissance r\'eelle du dividende de
      l'indice~; le capital n'est jamais vendu \\
    R\'eussite & Non mesur\'ee
      & Revenu d'au moins \SeuilReussite~\% de l'objectif chaque ann\'ee pendant \McRenteAns~ans \\
    R\'eserve & Aucune
      & Aucune, sauf dans la variante de la section~\ref{sec:risques-reserve} \\
    \bottomrule
  \end{tabularx}
\end{table}

Ce tableau montre que le mod\`ele de ce chapitre est bien plus indulgent que la projection sur un
point d\'ecisif~: il ignore l'imp\^ot et l'assurance maladie. Il juge l'objectif atteint d\`es que le
capital vaut \McKapitalSeuil~€, soit \McKapitalSeuilPartProzent~\% des \KapitalRequisReference~€ que le
chapitre~\ref{sec:capital} demande pour la r\'epartition \ReferenzStrategie. Il atteint donc
l'objectif beaucoup plus t\^ot~: \`a \McAgeObjectifMedian~ans pour la trajectoire m\'ediane du Monte
Carlo, contre \AlterYannBasis~ans dans la projection. \bk{Ne lis jamais un \^age ou un capital de ce
chapitre comme le tien~; ce sont ceux de la projection du chapitre~\ref{sec:accumulation}.} Ce que ce
mod\`ele mesure, c'est autre chose~: la probabilit\'e que le dividende d'un grand march\'e tienne
pendant toute une retraite, et la fa\c{c}on dont cette probabilit\'e change avec un choc ou une
r\'eserve. Pour cette question, l'imp\^ot joue peu, puisqu'il pr\'el\`eve \`a peu pr\`es la m\^eme part
d'un dividende haut ou bas.

Deux autres limites comptent. Le mod\`ele applique le rendement de l'ETF maison \`a tout le capital et
fait \'evoluer ton revenu comme le dividende de l'indice am\'ericain~; la r\'epartition
\ReferenzStrategie{} n'y entre que par le plan d'\'epargne, et ni l'ETF monde ni tes titres europ\'eens
n'y figurent. Il ne rejoue pas non plus la r\'ef\'erence en ETF avec retrait de \TauxRetrait~\%~: ce rejeu, et la
comparaison des deux strat\'egies sur les m\^emes ann\'ees de d\'epart, font l'objet de la
section~\ref{sec:retraite-quatre}.

\subsection{Rejouer l'histoire ann\'ee par ann\'ee}\label{sec:risques-rejeu}

Le rejeu pose une question simple~: si tu avais commenc\'e ton plan une ann\'ee donn\'ee du pass\'e,
comment l'histoire l'aurait-elle trait\'e~? La figure~\ref{fig:rejeu-annee-depart} r\'epond pour
chaque ann\'ee de d\'epart de la s\'erie.

\annahme{Pour chaque ann\'ee de d\'epart de \RueckspielPremier{} \`a \RueckspielDernier, ton plan
d'\'epargne \`a \ReferenzSparquote~\% re\c{c}oit les rendements r\'eels de l'indice dans l'ordre
historique jusqu'\`a l'objectif au sens du tableau~\ref{tab:projection-contre-simulation}. Le revenu
suit ensuite le dividende r\'eel de l'indice pendant \McRenteAns~ans~; une seule ann\'ee sous
\SeuilReussite~\% de l'objectif compte comme un \'echec, sans r\'eserve. Un d\'epart dont l'objectif
ou la retraite d\'epasse \ShillerFin{} n'est pas jug\'e.}

% Trois classes de points: une colonne jumelle de jahre_bis_ziel par issue (kapitel_7);
% les cellules vides sont lues comme nan et ecartees du nuage.
\linienfigur{fig_rueckspiel}[xlabel={Ann\'ee de d\'epart du plan d'\'epargne},
  ylabel={Ann\'ees d'\'epargne jusqu'\`a l'objectif},
  xticklabel style={/pgf/number format/set thousands separator={}},
  table/empty cells with={nan}, unbounded coords=discard,
  cycle list={{blue},{red},{gray}}, every axis plot post/.append style={only marks, mark=*,
    mark size=1.3pt},
  legend cell align=left, legend style={at={(0.5,-0.22)}, anchor=north, legend columns=3,
    /tikz/every even column/.append style={column sep=0.5cm}}]%
  {start}{ans_tenu,ans_echec,ans_non_juge}%
  {{Revenu tenu},{Revenu tomb\'e sous le seuil},{Retraite non observable jusqu'au bout}}%
  {Rejeu historique du plan d'\'epargne \`a \ReferenzSparquote~\% par ann\'ee de d\'epart~: nombre
  d'ann\'ees jusqu'\`a l'objectif du mod\`ele sans imp\^ot et issue de la retraite de \McRenteAns~ans
  (indice \IndiceActionsUs, \ShillerDebut--\ShillerFin).\label{fig:rejeu-annee-depart}}

Ce graphique montre que sur \RueckspielDebuts{} ann\'ees de d\'epart, \RueckspielAtteints{}
atteignent l'objectif du mod\`ele avant la fin des donn\'ees, en \RueckspielAnsMin{} \`a
\RueckspielAnsMax~ans d'\'epargne selon l'ann\'ee, avec une m\'ediane de \RueckspielAnsMedian~ans. Les
d\'eparts post\'erieurs \`a \RueckspielDernierAtteint{} n'ont pas encore eu le temps d'y arriver. Pour
\RueckspielJuges{} d\'eparts, jusqu'\`a celui de \RueckspielDernierJuge, l'histoire couvre aussi les
\McRenteAns~ans de retraite~: \RueckspielTenus{} tiennent leur revenu au-dessus du seuil, soit
\RueckspielTenusProzent~\%, et \RueckspielEchecs{} \'echouent. \bk{Les \RueckspielEchecs{} \'echecs
ont tous la m\^eme cause~: le revenu passe sous le seuil pendant l'une des
\RueckspielEchecsNombreEpisodes{} baisses de \RueckspielEchecsEpisodes{} de la
figure~\ref{fig:baisses-dividende-crises}.} Ce sont des retraites commenc\'ees avant ou pendant ces
baisses, dont le revenu n'avait pas encore assez progress\'e pour amortir le choc.

Dans ce mod\`ele, un krach pendant l'\'epargne ne fait qu'allonger la dur\'ee d'\'epargne, de
\RueckspielAnsMin{} \`a \RueckspielAnsMax~ans selon l'ann\'ee de d\'epart~; c'est la baisse du dividende
une fois \`a la retraite qui fait \'echouer le plan. Il faut lire ce pourcentage avec prudence. Les
\RueckspielJuges{} d\'eparts jug\'es ne sont pas autant d'exp\'eriences ind\'ependantes~: deux ann\'ees
voisines partagent presque toute leur histoire, et la p\'eriode \ShillerDebut--\ShillerFin{} ne
contient qu'une poign\'ee de retraites de \McRenteAns~ans sans chevauchement. Le tirage Monte Carlo
sert \`a d\'epasser cette limite.

\subsection{Des milliers d'histoires tir\'ees au hasard}\label{sec:risques-montecarlo}

Le Monte Carlo fabrique des histoires qui n'ont pas eu lieu mais qui auraient pu avoir lieu, en
recombinant des morceaux de l'histoire r\'eelle. La figure~\ref{fig:eventail-monte-carlo} montre
l'\'eventail du capital qui en r\'esulte.

\annahme{\McTrajectoires{} trajectoires de \McDureeAns~ans, chacune form\'ee de blocs de \McBloc~ann\'ees
cons\'ecutives tir\'es au hasard dans la s\'erie \ShillerDebut--\ShillerFin. Un bloc garde
l'encha\^inement r\'eel d'une crise et de sa reprise~; un tirage ann\'ee par ann\'ee le d\'etruirait.
Graine fixe, donc r\'esultats reproductibles. M\^emes r\`egles d'objectif et de r\'eussite que le rejeu
(tableau~\ref{tab:projection-contre-simulation}). Le capital n'est jamais retir\'e~: apr\`es
l'objectif, il continue de cro\^itre dividendes r\'einvestis, ce que tu ne feras pas en vivant de ces
dividendes.}

\linienfigur{fig_mc_faecher}[xlabel={Ann\'ee}, ylabel={Capital (millions d'euros de 2026)},
  xticklabel style={/pgf/number format/set thousands separator={}}, xmax=\AnneeHorizonFin,
  y filter/.expression={y/1e6}, ymin=0,
  cycle list={{blue!50},{blue, very thick},{blue!50},{red, dashed},{black, dashed}},
  legend pos=north west, legend cell align=left, legend style={fill=white}]%
  {jahr}{p10,p50,p90,seuil_mc,capital_requis}%
  {{\McCentileBas\textsuperscript{e} centile},{M\'ediane},{\McCentileHaut\textsuperscript{e} centile},
   {Seuil du Monte Carlo (sans imp\^ot)},{Capital requis du chapitre~\ref{sec:capital} (avec imp\^ot)}}%
  {\'Eventail Monte Carlo du capital du plan d'\'epargne \`a \ReferenzSparquote~\% jusqu'en
  \AnneeHorizonFin~: m\'ediane et centiles des \McTrajectoires{} trajectoires, face au seuil du mod\`ele
  sans imp\^ot et au capital requis de la r\'epartition \ReferenzStrategie{} (euros de
  2026).\label{fig:eventail-monte-carlo}}

Ce graphique montre qu'en \AnneeHorizonFin{} le capital m\'edian atteint \McCapitalHorizonMedian~€,
mais qu'une trajectoire sur dix reste sous \McCapitalHorizonBas~€ et qu'une sur dix d\'epasse
\McCapitalHorizonHaut~€~: \McEventailRapport{} fois plus, pour la m\^eme \'epargne. Seul le hasard des
rendements tir\'es s\'epare ces trajectoires. La m\'ediane franchit le seuil sans imp\^ot en
\McAnneeObjectifMedian, soit \`a \McAgeObjectifMedian~ans, et huit trajectoires sur dix le franchissent
entre \McAnneeObjectifBas{} et \McAnneeObjectifHaut. La ligne du capital requis avec imp\^ot, bien plus
haute, rappelle pourquoi ces dates ne sont pas les tiennes. Apr\`es l'objectif, les courbes montent
encore parce que le mod\`ele ne retire rien~; seule la partie avant l'objectif d\'ecrit un capital que
tu poss\'ederais.

Sur ces trajectoires, \bk{le revenu tient au-dessus de \SeuilReussite~\% de l'objectif pendant
\McRenteAns~ans dans \ErfolgsquoteBasis~\% des cas, sans r\'eserve}, un r\'esultat proche des
\RueckspielTenusProzent~\% du rejeu. Ce taux compte aussi des trajectoires qui n'atteignent l'objectif
qu'apr\`es \AnneeHorizonFin, parfois des d\'ecennies plus tard~: si l'on exige l'objectif dans
l'horizon de ton plan, il tombe \`a \McErfolgDansHorizon~\%, \McErfolgEcartHorizon~points de moins. La
dur\'ee de la retraite joue peu~: la r\'eussite vaut \McErfolgCourt~\% sur \McRenteAnsCourt~ans et
\McErfolgLong~\% sur \McRenteAnsLong~ans, ce qui signifie que la plupart des \'echecs surviennent dans
les \McRenteAnsCourt{} premi\`eres ann\'ees. Le Monte Carlo ne tire que dans l'histoire am\'ericaine,
comme le rejeu~: il ne dit rien d'un march\'e europ\'een, dont le dividende peut \^etre plus
irr\'egulier.

\subsection{Le choc de la premi\`ere ann\'ee de retraite}\label{sec:risques-sequence}

Le rejeu a montr\'e que les \'echecs viennent des baisses qui frappent pendant la retraite, avant que
le revenu ait pu progresser. La figure~\ref{fig:reussite-choc-dividende} provoque d\'elib\'er\'ement
une telle baisse, ou une hausse, la premi\`ere ann\'ee.

\annahme{\HypotheseDivSchock. Grille de $-$\ChocGrilleMax{} \`a $+$\ChocGrilleMax~points par pas de
\ChocGrillePas~points, sur les m\^emes trajectoires que la figure~\ref{fig:eventail-monte-carlo}, pour
des retraites de \McRenteAnsCourt, \McRenteAns{} et \McRenteAnsLong~ans, sans r\'eserve.}

\linienfigur{fig_mc_erfolg}[xlabel={Choc sur la croissance r\'eelle du dividende la premi\`ere ann\'ee
  (points de pourcentage)}, ylabel={Probabilit\'e de r\'eussite (\%)},
  x filter/.expression={x*1e2}, y filter/.expression={y*1e2}, ymin=0,
  every axis plot post/.append style={mark=*, mark size=1.2pt},
  legend pos=south east, legend cell align=left, legend style={fill=white}]%
  {div_schock}{erfolg_30,erfolg_40,erfolg_50}%
  {{Retraite de \McRenteAnsCourt~ans},{Retraite de \McRenteAns~ans},{Retraite de \McRenteAnsLong~ans}}%
  {Probabilit\'e que le revenu reste au-dessus de \SeuilReussite~\% de l'objectif pendant toute la
  retraite, selon un choc sur la croissance r\'eelle du dividende la premi\`ere ann\'ee de retraite
  (Monte Carlo sans imp\^ot, \McTrajectoires{} trajectoires).\label{fig:reussite-choc-dividende}}

Ce graphique montre que la r\'eussite d\'epend fortement de la premi\`ere ann\'ee. Sans choc, elle vaut
\ErfolgsquoteBasis~\% pour une retraite de \McRenteAns~ans. Un choc de $-$10~points la ram\`ene \`a
\McErfolgChocMoinsDix~\%, un choc de $-$20~points \`a \McErfolgChocMoinsVingt~\%, et un choc de
$-$\ChocGrilleMax~points \`a \McErfolgChocMin~\%. \`A l'inverse, une premi\`ere ann\'ee favorable de
$+$\ChocGrilleMax~points la porte \`a \McErfolgChocMax~\%. \bk{Un choc de $-$20~points, moins que la
baisse de la crise financi\`ere (\EinbruchRecent~\%, \'etal\'ee de \EinbruchRecentDebut{} \`a
\EinbruchRecentFin), r\'eduit d\'ej\`a tes chances de plus de moiti\'e s'il tombe la premi\`ere ann\'ee
de ta retraite.} La raison est m\'ecanique~: la r\`egle tol\`ere une baisse de revenu jusqu'au seuil de
\SeuilReussite~\%, et un tel choc consomme d'un coup presque toute cette marge. Les trois courbes
restent proches, ce qui confirme que la dur\'ee de la retraite compte moins que son d\'ebut.

Ce choc sert aussi \`a approcher ce que l'histoire am\'ericaine ne contient pas~: une coupe plus
brutale sur un portefeuille plus concentr\'e que l'indice, comme celles que peut subir un ETF maison de
\NombreTitres{} titres. Le mod\`ele ne le teste qu'\`a la premi\`ere ann\'ee de retraite~; il ne dit
pas ce que le m\^eme choc ferait dix ans plus tard.

\subsection{La r\'eserve de liquidit\'es}\label{sec:risques-reserve}

Le rem\`ede classique au risque du d\'ebut de retraite est une r\'eserve en liquidit\'es, dans laquelle
tu puises quand le dividende baisse au lieu de r\'eduire tes d\'epenses.

\annahme{R\'eserve de \PufferJahreMit~ans d'objectif, soit \PufferMontantMit~€, d\'etenue en
liquidit\'es au d\'epart en retraite, sans int\'er\^et ni imp\^ot. Elle ne comble que ce qui manque
pour remonter au seuil de \SeuilReussite~\% de l'objectif, jamais l'\'ecart restant jusqu'\`a
l'objectif, et n'est jamais reconstitu\'ee.}

Sur les m\^emes trajectoires, cette r\'eserve fait passer la r\'eussite de \ErfolgsquoteBasis~\% \`a
\ErfolgsquoteMitPuffer~\%. C'est l'effet le plus fort de ce chapitre~: \PufferJahreMit~ann\'ees de
d\'epenses mises de c\^ot\'e suffisent \`a traverser la plupart des baisses qui font \'echouer le plan
sans r\'eserve. Ce confort a un prix. Les \PufferMontantMit~€ repr\'esentent
\PufferPartCapitalProzent~\% du capital requis de la r\'epartition \ReferenzStrategie{} et ne sont pas
compt\'es dans les capitaux du chapitre~\ref{sec:capital}~: il faut les \'epargner en plus, et ils ne
rapportent rien. La simulation suppose en outre que tu acceptes de vivre avec \SeuilReussite~\% de ta
cible pendant une baisse, la r\'eserve servant seulement \`a ne pas descendre plus bas~; si tu la
puisais pour garder ton niveau de vie entier, elle s'\'epuiserait plus vite. La
section~\ref{sec:retraite-reserve} montre l'effet de la r\'eserve selon sa taille.

\subsection{L'inflation de 2022 et la rente non index\'ee}\label{sec:risques-inflation}

La cible de ce plan est fix\'ee en euros de 2026~: chaque ann\'ee, ton revenu doit suivre les prix. Un
revenu fix\'e une fois pour toutes en euros courants, comme une rente priv\'ee non revaloris\'ee, ne le
fait pas. La figure~\ref{fig:inflation-rente-fixe} en montre le co\^ut sur les derni\`eres ann\'ees.

\annahme{Inflation annuelle allemande publi\'ee par Destatis de \InflationSerieDebut{} \`a
\InflationSerieFin\QH{inflation_destatis_2019_2025}. Pour \InflationPicAnnee, Destatis annon\c{c}ait
d'abord \InflationPicProvisoire~\%\QH{inflation_2022_de}, et ses communiqu\'es ult\'erieurs citent
\InflationPicRevisee~\%~; la figure retient la s\'erie r\'evis\'ee. Rente fixe de \ZielMonat~€ par mois
\`a partir de \InflationSerieDebut, jamais revaloris\'ee, compar\'ee \`a un revenu qui suit les prix.}

\linienfigur{fig_inflation_2022}[xlabel={Ann\'ee},
  ylabel={Pouvoir d'achat mensuel (euros de \InflationSerieDebut)},
  xticklabel style={/pgf/number format/set thousands separator={}}, scaled y ticks=false,
  width=0.95\linewidth, xtick=data,
  every axis plot post/.append style={mark=*, mark size=1.5pt},
  legend pos=south west, legend cell align=left, legend style={fill=white}]%
  {jahr}{indexiert,nicht_indexiert}%
  {{Revenu qui suit les prix},{Rente fixe non index\'ee}}%
  {Pouvoir d'achat d'un revenu de \ZielMonat~€ par mois \`a partir de \InflationSerieDebut, selon qu'il
  suit les prix ou reste fixe en euros courants (inflation allemande publi\'ee par Destatis, euros de
  \InflationSerieDebut).\label{fig:inflation-rente-fixe}}

Ce graphique montre qu'une rente fixe de \ZielMonat~€ ne valait plus que \RenteNonIndexeeFin~€ en
\InflationSerieFin, soit \RentePerteMois~€ par mois ou \RentePerteProzent~\% de pouvoir d'achat en
moins, apr\`es une hausse des prix de \PrixHausseCumulee~\%. Plus de la moiti\'e de cette hausse est
venue de \InflationPicAnnee{} et de 2023. Une inflation calme finit par co\^uter autant~: \`a
\InflationZiel~\% par an, l'objectif de la Banque centrale europ\'eenne\QH{inflation_ziel}, une rente
fixe perd la moiti\'e de sa valeur en \RenteMoitieAns~ans, et apr\`es \McRenteAns~ans de retraite
\ZielMonat~€ n'en valent plus que \RenteApresRetraite~€.

Les dividendes ne sont index\'es par aucun contrat. Sur longue p\'eriode, celui de l'indice am\'ericain a
battu l'inflation (section~\ref{sec:risques-histoire}), et il a encore progress\'e en termes r\'eels en
\ShillerFin. Mais pendant la d\'ecennie d'inflation de \EinbruchInflationDebut{} \`a
\EinbruchInflationFin, il a perdu \EinbruchInflation~\% de son pouvoir d'achat, et des retraites
commenc\'ees au d\'ebut de cette d\'ecennie ont \'echou\'e dans le rejeu. Toute part fixe en euros dans
tes revenus futurs subit l'inflation en entier~; la part en dividendes la subit quand les entreprises
ne suivent pas les prix.

\bk{Retiens de ce chapitre que le dividende d'un grand march\'e a d\'ej\`a perdu jusqu'\`a
\GroessterEinbruch~\% de sa valeur r\'eelle et est rest\'e jusqu'\`a \BaisseLongueAns~ans sous son
sommet, que le vrai danger est une baisse au d\'ebut de la retraite, et qu'une r\'eserve de
\PufferJahreMit~ans de d\'epenses en prot\`ege dans la plupart des cas.} Dans le mod\`ele sans imp\^ot de
ce chapitre, ton revenu tient \McRenteAns~ans dans \ErfolgsquoteBasis~\% des cas sans r\'eserve et dans
\ErfolgsquoteMitPuffer~\% avec~; ces taux mesurent la solidit\'e du dividende, pas la date de ton
d\'epart. S'y ajoute l'hypoth\`ese favorable d'un dividende qui cro\^it bien plus vite que par le
pass\'e. Le chapitre~\ref{sec:retraite} passe \`a la retraite elle-m\^eme~: la p\'eriode avant la rente l\'egale, le
montant de cette rente et la taille de la r\'eserve.

\Yannbox{\AnneeYannBasis}{Dans la projection, Yann atteint son objectif en \AnneeYannBasis, \`a
\AlterYannBasis~ans, avec \BruttoNoetigMois~€ de dividendes bruts par mois. Si ses dividendes
subissaient la pire baisse de l'histoire de l'\IndiceActionsUs, \GroessterEinbruch~\% comme de
\GroessterEinbruchDebut{} \`a \GroessterEinbruchJahr, ils tomberaient \`a \YannBrutApresPireBaisse~€
bruts par mois, et l'indice a mis \GroessterEinbruchSousSommetAns~ans \`a retrouver son sommet. Dans la
simulation sans imp\^ot, son revenu tient \McRenteAns~ans dans \ErfolgsquoteBasis~\% des cas, et dans
\ErfolgsquoteMitPuffer~\% s'il garde \PufferMontantMit~€ de r\'eserve, \PufferJahreMit~ans de d\'epenses
\`a \'epargner en plus. Si une partie de ses revenus \'etait fixe en euros, une inflation comme celle de
\InflationSerieDebut{} \`a \InflationSerieFin{} lui en retirerait \RentePerteProzent~\%.}
